% !TEX root = ../main.tex
% Convolutional encoder in the bb_theme figure style (palette, wires and vector cells
% follow bb_theme/latex/vae-canonical/vae.tex). Illustrative sizes.
% Use inside a figure environment:  % !TEX root = ../main.tex
% Convolutional encoder in the bb_theme figure style (palette, wires and vector cells
% follow bb_theme/latex/vae-canonical/vae.tex). Illustrative sizes.
% Use inside a figure environment:  % !TEX root = ../main.tex
% Convolutional encoder in the bb_theme figure style (palette, wires and vector cells
% follow bb_theme/latex/vae-canonical/vae.tex). Illustrative sizes.
% Use inside a figure environment:  % !TEX root = ../main.tex
% Convolutional encoder in the bb_theme figure style (palette, wires and vector cells
% follow bb_theme/latex/vae-canonical/vae.tex). Illustrative sizes.
% Use inside a figure environment:  \input{figures/cnn.tex}
% Needs: tikz (arrows.meta), graphicx, amsmath, amssymb.
\usetikzlibrary{positioning}
\providecommand{\cnnimg}{figures/cnn_input.png}
\definecolor{vaeInk}{HTML}{0F1B2D}
\definecolor{vaeHatch}{HTML}{2F3FD0}
\definecolor{vaeSignal}{HTML}{1D5FD1}
\definecolor{vaeMuted}{HTML}{56677F}
\definecolor{vaeFill}{HTML}{EEF3FB}
\begin{tikzpicture}[
    font=\small, line cap=round, line join=round,
    wire/.style={draw=vaeInk, line width=0.6pt, -{Stealth[length=4.5pt, width=3.5pt]}},
    proj/.style={draw=vaeHatch, line width=0.4pt, dash pattern=on 2.5pt off 1.8pt},
    net/.style={draw=vaeInk, line width=0.7pt, fill=vaeFill},
    frame/.style={draw=vaeInk, line width=0.6pt},
    vcap/.style={font=\footnotesize, text=vaeMuted, align=center},
  ]
  \def\cs{0.17}   % vector cell size
  \def\d{0.11}    % offset between stacked channels

  % ---------- input depth image ----------
  \node[inner sep=0] (din) at (1.0,0) {\includegraphics[width=2.0cm, height=1.5cm]{\cnnimg}};
  \draw[frame] (din.south west) rectangle (din.north east);
  \node[above=2pt of din.north] {$D_t$};
  \node[vcap, below=3pt of din.south] {$1{\times}H{\times}W$};
  % kernel window on the input
  \draw[vaeHatch, line width=0.8pt, fill=vaeHatch!15] (0.35,0.3) rectangle ++(0.34,0.34);

  % ---------- feature maps: back to front ----------
  % map 1
  \foreach \i/\s in {3/12,2/22,1/32,0/45} {
    \draw[frame, fill=vaeSignal!\s!white] (3.35+\d*\i,-0.62+\d*\i) rectangle ++(1.5,1.15);
  }
  \node[vcap, below=3pt] at (3.35+0.75+0.17,-0.62) {conv, ReLU\\$C_1{\times}\frac{H}{2}{\times}\frac{W}{2}$};
  % map 2
  \foreach \i/\s in {4/12,3/22,2/32,1/45,0/60} {
    \draw[frame, fill=vaeSignal!\s!white] (5.75+\d*\i,-0.4+\d*\i) rectangle ++(1.0,0.78);
  }
  \node[vcap, below=3pt] at (5.75+0.5+0.22,-0.4) {conv, ReLU\\$C_2{\times}\frac{H}{4}{\times}\frac{W}{4}$};
  % map 3
  \foreach \i/\s in {5/12,4/22,3/32,2/45,1/60,0/75} {
    \draw[frame, fill=vaeSignal!\s!white] (7.65+\d*\i,-0.26) ++(0,\d*\i) rectangle ++(0.7,0.54);
  }
  \node[vcap, below=3pt] at (7.65+0.35+0.28,-0.26) {conv, ReLU\\$C_3{\times}\frac{H}{8}{\times}\frac{W}{8}$};

  % ---------- kernel projection onto the first map ----------
  \draw[vaeHatch, line width=0.8pt, fill=vaeHatch!15] (3.35+0.5,-0.62+0.55) rectangle ++(0.13,0.13);
  \draw[proj] (0.69,0.64) -- (3.85+0.0,-0.62+0.68);
  \draw[proj] (0.69,0.30) -- (3.85+0.0,-0.62+0.55);
  \node[font=\footnotesize, text=vaeHatch, anchor=south] at (2.1,1.05) {kernel};

  % ---------- wires between stages ----------
  \draw[wire] (din.east) ++(0,-0.0) -- (3.2,0) node[midway, below=1pt, vcap] {stride 2};
  \draw[wire] (5.05+0.0,0) -- (5.95-0.2,0);
  \draw[wire] (7.05+0.2,0) -- (7.75-0.1,0);

  % ---------- flatten: grid of cells ----------
  \foreach \i/\s in {0/25,1/60,2/40,3/80,4/15,5/50,6/70,7/30,8/55} {
    \fill[vaeSignal!\s!white] (9.45,-0.7+\i*\cs) rectangle ++(\cs,\cs);
  }
  \draw[vaeInk, line width=0.5pt, step=\cs, shift={(9.45,-0.7)}] (0,0) grid (\cs,9*\cs);
  \node[vcap, below=3pt] at (9.45+\cs/2,-0.7) {flatten};
  \draw[wire] (8.5,0) -- (9.4,0);

  % ---------- linear layer ----------
  \draw[net] (10.0,0.5) -- (11.0,0.3) -- (11.0,-0.3) -- (10.0,-0.5) -- cycle;
  \node[font=\footnotesize] at (10.5,0) {Linear};
  \draw[wire] (9.45+\cs,0) -- (10.0,0);

  % ---------- feature vector z ----------
  \foreach \i/\s in {0/55,1/20,2/70,3/35,4/45,5/10} {
    \fill[vaeSignal!\s!white] (11.55,-3*\cs+\i*\cs) rectangle ++(\cs,\cs);
  }
  \draw[vaeSignal, line width=0.7pt, step=\cs, shift={(11.55,-3*\cs)}] (0,0) grid (\cs,6*\cs);
  \draw[wire] (11.0,0) -- (11.55,0);
  \node[anchor=south, inner sep=1pt] at (11.55+\cs/2,3*\cs+0.08) {$\mathbf{z}\in\mathbb{R}^{d}$};
\end{tikzpicture}

% Needs: tikz (arrows.meta), graphicx, amsmath, amssymb.
\usetikzlibrary{positioning}
\providecommand{\cnnimg}{figures/cnn_input.png}
\definecolor{vaeInk}{HTML}{0F1B2D}
\definecolor{vaeHatch}{HTML}{2F3FD0}
\definecolor{vaeSignal}{HTML}{1D5FD1}
\definecolor{vaeMuted}{HTML}{56677F}
\definecolor{vaeFill}{HTML}{EEF3FB}
\begin{tikzpicture}[
    font=\small, line cap=round, line join=round,
    wire/.style={draw=vaeInk, line width=0.6pt, -{Stealth[length=4.5pt, width=3.5pt]}},
    proj/.style={draw=vaeHatch, line width=0.4pt, dash pattern=on 2.5pt off 1.8pt},
    net/.style={draw=vaeInk, line width=0.7pt, fill=vaeFill},
    frame/.style={draw=vaeInk, line width=0.6pt},
    vcap/.style={font=\footnotesize, text=vaeMuted, align=center},
  ]
  \def\cs{0.17}   % vector cell size
  \def\d{0.11}    % offset between stacked channels

  % ---------- input depth image ----------
  \node[inner sep=0] (din) at (1.0,0) {\includegraphics[width=2.0cm, height=1.5cm]{\cnnimg}};
  \draw[frame] (din.south west) rectangle (din.north east);
  \node[above=2pt of din.north] {$D_t$};
  \node[vcap, below=3pt of din.south] {$1{\times}H{\times}W$};
  % kernel window on the input
  \draw[vaeHatch, line width=0.8pt, fill=vaeHatch!15] (0.35,0.3) rectangle ++(0.34,0.34);

  % ---------- feature maps: back to front ----------
  % map 1
  \foreach \i/\s in {3/12,2/22,1/32,0/45} {
    \draw[frame, fill=vaeSignal!\s!white] (3.35+\d*\i,-0.62+\d*\i) rectangle ++(1.5,1.15);
  }
  \node[vcap, below=3pt] at (3.35+0.75+0.17,-0.62) {conv, ReLU\\$C_1{\times}\frac{H}{2}{\times}\frac{W}{2}$};
  % map 2
  \foreach \i/\s in {4/12,3/22,2/32,1/45,0/60} {
    \draw[frame, fill=vaeSignal!\s!white] (5.75+\d*\i,-0.4+\d*\i) rectangle ++(1.0,0.78);
  }
  \node[vcap, below=3pt] at (5.75+0.5+0.22,-0.4) {conv, ReLU\\$C_2{\times}\frac{H}{4}{\times}\frac{W}{4}$};
  % map 3
  \foreach \i/\s in {5/12,4/22,3/32,2/45,1/60,0/75} {
    \draw[frame, fill=vaeSignal!\s!white] (7.65+\d*\i,-0.26) ++(0,\d*\i) rectangle ++(0.7,0.54);
  }
  \node[vcap, below=3pt] at (7.65+0.35+0.28,-0.26) {conv, ReLU\\$C_3{\times}\frac{H}{8}{\times}\frac{W}{8}$};

  % ---------- kernel projection onto the first map ----------
  \draw[vaeHatch, line width=0.8pt, fill=vaeHatch!15] (3.35+0.5,-0.62+0.55) rectangle ++(0.13,0.13);
  \draw[proj] (0.69,0.64) -- (3.85+0.0,-0.62+0.68);
  \draw[proj] (0.69,0.30) -- (3.85+0.0,-0.62+0.55);
  \node[font=\footnotesize, text=vaeHatch, anchor=south] at (2.1,1.05) {kernel};

  % ---------- wires between stages ----------
  \draw[wire] (din.east) ++(0,-0.0) -- (3.2,0) node[midway, below=1pt, vcap] {stride 2};
  \draw[wire] (5.05+0.0,0) -- (5.95-0.2,0);
  \draw[wire] (7.05+0.2,0) -- (7.75-0.1,0);

  % ---------- flatten: grid of cells ----------
  \foreach \i/\s in {0/25,1/60,2/40,3/80,4/15,5/50,6/70,7/30,8/55} {
    \fill[vaeSignal!\s!white] (9.45,-0.7+\i*\cs) rectangle ++(\cs,\cs);
  }
  \draw[vaeInk, line width=0.5pt, step=\cs, shift={(9.45,-0.7)}] (0,0) grid (\cs,9*\cs);
  \node[vcap, below=3pt] at (9.45+\cs/2,-0.7) {flatten};
  \draw[wire] (8.5,0) -- (9.4,0);

  % ---------- linear layer ----------
  \draw[net] (10.0,0.5) -- (11.0,0.3) -- (11.0,-0.3) -- (10.0,-0.5) -- cycle;
  \node[font=\footnotesize] at (10.5,0) {Linear};
  \draw[wire] (9.45+\cs,0) -- (10.0,0);

  % ---------- feature vector z ----------
  \foreach \i/\s in {0/55,1/20,2/70,3/35,4/45,5/10} {
    \fill[vaeSignal!\s!white] (11.55,-3*\cs+\i*\cs) rectangle ++(\cs,\cs);
  }
  \draw[vaeSignal, line width=0.7pt, step=\cs, shift={(11.55,-3*\cs)}] (0,0) grid (\cs,6*\cs);
  \draw[wire] (11.0,0) -- (11.55,0);
  \node[anchor=south, inner sep=1pt] at (11.55+\cs/2,3*\cs+0.08) {$\mathbf{z}\in\mathbb{R}^{d}$};
\end{tikzpicture}

% Needs: tikz (arrows.meta), graphicx, amsmath, amssymb.
\usetikzlibrary{positioning}
\providecommand{\cnnimg}{figures/cnn_input.png}
\definecolor{vaeInk}{HTML}{0F1B2D}
\definecolor{vaeHatch}{HTML}{2F3FD0}
\definecolor{vaeSignal}{HTML}{1D5FD1}
\definecolor{vaeMuted}{HTML}{56677F}
\definecolor{vaeFill}{HTML}{EEF3FB}
\begin{tikzpicture}[
    font=\small, line cap=round, line join=round,
    wire/.style={draw=vaeInk, line width=0.6pt, -{Stealth[length=4.5pt, width=3.5pt]}},
    proj/.style={draw=vaeHatch, line width=0.4pt, dash pattern=on 2.5pt off 1.8pt},
    net/.style={draw=vaeInk, line width=0.7pt, fill=vaeFill},
    frame/.style={draw=vaeInk, line width=0.6pt},
    vcap/.style={font=\footnotesize, text=vaeMuted, align=center},
  ]
  \def\cs{0.17}   % vector cell size
  \def\d{0.11}    % offset between stacked channels

  % ---------- input depth image ----------
  \node[inner sep=0] (din) at (1.0,0) {\includegraphics[width=2.0cm, height=1.5cm]{\cnnimg}};
  \draw[frame] (din.south west) rectangle (din.north east);
  \node[above=2pt of din.north] {$D_t$};
  \node[vcap, below=3pt of din.south] {$1{\times}H{\times}W$};
  % kernel window on the input
  \draw[vaeHatch, line width=0.8pt, fill=vaeHatch!15] (0.35,0.3) rectangle ++(0.34,0.34);

  % ---------- feature maps: back to front ----------
  % map 1
  \foreach \i/\s in {3/12,2/22,1/32,0/45} {
    \draw[frame, fill=vaeSignal!\s!white] (3.35+\d*\i,-0.62+\d*\i) rectangle ++(1.5,1.15);
  }
  \node[vcap, below=3pt] at (3.35+0.75+0.17,-0.62) {conv, ReLU\\$C_1{\times}\frac{H}{2}{\times}\frac{W}{2}$};
  % map 2
  \foreach \i/\s in {4/12,3/22,2/32,1/45,0/60} {
    \draw[frame, fill=vaeSignal!\s!white] (5.75+\d*\i,-0.4+\d*\i) rectangle ++(1.0,0.78);
  }
  \node[vcap, below=3pt] at (5.75+0.5+0.22,-0.4) {conv, ReLU\\$C_2{\times}\frac{H}{4}{\times}\frac{W}{4}$};
  % map 3
  \foreach \i/\s in {5/12,4/22,3/32,2/45,1/60,0/75} {
    \draw[frame, fill=vaeSignal!\s!white] (7.65+\d*\i,-0.26) ++(0,\d*\i) rectangle ++(0.7,0.54);
  }
  \node[vcap, below=3pt] at (7.65+0.35+0.28,-0.26) {conv, ReLU\\$C_3{\times}\frac{H}{8}{\times}\frac{W}{8}$};

  % ---------- kernel projection onto the first map ----------
  \draw[vaeHatch, line width=0.8pt, fill=vaeHatch!15] (3.35+0.5,-0.62+0.55) rectangle ++(0.13,0.13);
  \draw[proj] (0.69,0.64) -- (3.85+0.0,-0.62+0.68);
  \draw[proj] (0.69,0.30) -- (3.85+0.0,-0.62+0.55);
  \node[font=\footnotesize, text=vaeHatch, anchor=south] at (2.1,1.05) {kernel};

  % ---------- wires between stages ----------
  \draw[wire] (din.east) ++(0,-0.0) -- (3.2,0) node[midway, below=1pt, vcap] {stride 2};
  \draw[wire] (5.05+0.0,0) -- (5.95-0.2,0);
  \draw[wire] (7.05+0.2,0) -- (7.75-0.1,0);

  % ---------- flatten: grid of cells ----------
  \foreach \i/\s in {0/25,1/60,2/40,3/80,4/15,5/50,6/70,7/30,8/55} {
    \fill[vaeSignal!\s!white] (9.45,-0.7+\i*\cs) rectangle ++(\cs,\cs);
  }
  \draw[vaeInk, line width=0.5pt, step=\cs, shift={(9.45,-0.7)}] (0,0) grid (\cs,9*\cs);
  \node[vcap, below=3pt] at (9.45+\cs/2,-0.7) {flatten};
  \draw[wire] (8.5,0) -- (9.4,0);

  % ---------- linear layer ----------
  \draw[net] (10.0,0.5) -- (11.0,0.3) -- (11.0,-0.3) -- (10.0,-0.5) -- cycle;
  \node[font=\footnotesize] at (10.5,0) {Linear};
  \draw[wire] (9.45+\cs,0) -- (10.0,0);

  % ---------- feature vector z ----------
  \foreach \i/\s in {0/55,1/20,2/70,3/35,4/45,5/10} {
    \fill[vaeSignal!\s!white] (11.55,-3*\cs+\i*\cs) rectangle ++(\cs,\cs);
  }
  \draw[vaeSignal, line width=0.7pt, step=\cs, shift={(11.55,-3*\cs)}] (0,0) grid (\cs,6*\cs);
  \draw[wire] (11.0,0) -- (11.55,0);
  \node[anchor=south, inner sep=1pt] at (11.55+\cs/2,3*\cs+0.08) {$\mathbf{z}\in\mathbb{R}^{d}$};
\end{tikzpicture}

% Needs: tikz (arrows.meta), graphicx, amsmath, amssymb.
\usetikzlibrary{positioning}
\providecommand{\cnnimg}{figures/cnn_input.png}
\definecolor{vaeInk}{HTML}{0F1B2D}
\definecolor{vaeHatch}{HTML}{2F3FD0}
\definecolor{vaeSignal}{HTML}{1D5FD1}
\definecolor{vaeMuted}{HTML}{56677F}
\definecolor{vaeFill}{HTML}{EEF3FB}
\begin{tikzpicture}[
    font=\small, line cap=round, line join=round,
    wire/.style={draw=vaeInk, line width=0.6pt, -{Stealth[length=4.5pt, width=3.5pt]}},
    proj/.style={draw=vaeHatch, line width=0.4pt, dash pattern=on 2.5pt off 1.8pt},
    net/.style={draw=vaeInk, line width=0.7pt, fill=vaeFill},
    frame/.style={draw=vaeInk, line width=0.6pt},
    vcap/.style={font=\footnotesize, text=vaeMuted, align=center},
  ]
  \def\cs{0.17}   % vector cell size
  \def\d{0.11}    % offset between stacked channels

  % ---------- input depth image ----------
  \node[inner sep=0] (din) at (1.0,0) {\includegraphics[width=2.0cm, height=1.5cm]{\cnnimg}};
  \draw[frame] (din.south west) rectangle (din.north east);
  \node[above=2pt of din.north] {$D_t$};
  \node[vcap, below=3pt of din.south] {$1{\times}H{\times}W$};
  % kernel window on the input
  \draw[vaeHatch, line width=0.8pt, fill=vaeHatch!15] (0.35,0.3) rectangle ++(0.34,0.34);

  % ---------- feature maps: back to front ----------
  % map 1
  \foreach \i/\s in {3/12,2/22,1/32,0/45} {
    \draw[frame, fill=vaeSignal!\s!white] (3.35+\d*\i,-0.62+\d*\i) rectangle ++(1.5,1.15);
  }
  \node[vcap, below=3pt] at (3.35+0.75+0.17,-0.62) {conv, ReLU\\$C_1{\times}\frac{H}{2}{\times}\frac{W}{2}$};
  % map 2
  \foreach \i/\s in {4/12,3/22,2/32,1/45,0/60} {
    \draw[frame, fill=vaeSignal!\s!white] (5.75+\d*\i,-0.4+\d*\i) rectangle ++(1.0,0.78);
  }
  \node[vcap, below=3pt] at (5.75+0.5+0.22,-0.4) {conv, ReLU\\$C_2{\times}\frac{H}{4}{\times}\frac{W}{4}$};
  % map 3
  \foreach \i/\s in {5/12,4/22,3/32,2/45,1/60,0/75} {
    \draw[frame, fill=vaeSignal!\s!white] (7.65+\d*\i,-0.26) ++(0,\d*\i) rectangle ++(0.7,0.54);
  }
  \node[vcap, below=3pt] at (7.65+0.35+0.28,-0.26) {conv, ReLU\\$C_3{\times}\frac{H}{8}{\times}\frac{W}{8}$};

  % ---------- kernel projection onto the first map ----------
  \draw[vaeHatch, line width=0.8pt, fill=vaeHatch!15] (3.35+0.5,-0.62+0.55) rectangle ++(0.13,0.13);
  \draw[proj] (0.69,0.64) -- (3.85+0.0,-0.62+0.68);
  \draw[proj] (0.69,0.30) -- (3.85+0.0,-0.62+0.55);
  \node[font=\footnotesize, text=vaeHatch, anchor=south] at (2.1,1.05) {kernel};

  % ---------- wires between stages ----------
  \draw[wire] (din.east) ++(0,-0.0) -- (3.2,0) node[midway, below=1pt, vcap] {stride 2};
  \draw[wire] (5.05+0.0,0) -- (5.95-0.2,0);
  \draw[wire] (7.05+0.2,0) -- (7.75-0.1,0);

  % ---------- flatten: grid of cells ----------
  \foreach \i/\s in {0/25,1/60,2/40,3/80,4/15,5/50,6/70,7/30,8/55} {
    \fill[vaeSignal!\s!white] (9.45,-0.7+\i*\cs) rectangle ++(\cs,\cs);
  }
  \draw[vaeInk, line width=0.5pt, step=\cs, shift={(9.45,-0.7)}] (0,0) grid (\cs,9*\cs);
  \node[vcap, below=3pt] at (9.45+\cs/2,-0.7) {flatten};
  \draw[wire] (8.5,0) -- (9.4,0);

  % ---------- linear layer ----------
  \draw[net] (10.0,0.5) -- (11.0,0.3) -- (11.0,-0.3) -- (10.0,-0.5) -- cycle;
  \node[font=\footnotesize] at (10.5,0) {Linear};
  \draw[wire] (9.45+\cs,0) -- (10.0,0);

  % ---------- feature vector z ----------
  \foreach \i/\s in {0/55,1/20,2/70,3/35,4/45,5/10} {
    \fill[vaeSignal!\s!white] (11.55,-3*\cs+\i*\cs) rectangle ++(\cs,\cs);
  }
  \draw[vaeSignal, line width=0.7pt, step=\cs, shift={(11.55,-3*\cs)}] (0,0) grid (\cs,6*\cs);
  \draw[wire] (11.0,0) -- (11.55,0);
  \node[anchor=south, inner sep=1pt] at (11.55+\cs/2,3*\cs+0.08) {$\mathbf{z}\in\mathbb{R}^{d}$};
\end{tikzpicture}
